\section{Reduction to a strict exterior with a rest end}\label{four:sec:end}

This section reduces the invariant-mass assertion to an energy assertion
for one-ended data with compactly supported second fundamental form.
The replacement takes place arbitrarily far along the selected end.
It preserves the dominant energy condition after a correction of
vanishing energy cost. A proper map controls all enclosing cuts, including
cuts that enter the replacement region.

The endpoint used in this section is Theorem~\ref{four:thm:prepared-energy},
proved below: for the strict one-ended data described in
Lemma~\ref{four:lem:strictification},
\begin{equation}\label{four:end:prepared-interface}
 E_g\ge \frac12\left(\frac{A_*(g)}{\omega}\right)^{2/3}.
\end{equation}
Here and throughout this section, \(A_*(g)\) is the smooth-cut infimum
of Definition~\ref{four:def:cuts}, with the entire intrinsic frontier counted.
The reduction to \eqref{four:end:prepared-interface} does not use maximality.

\subsection{The exact cut class and removal of other ends}

\begin{lemma}[Cut reduction and positivity]\label{four:lem:cut-reduction}
Under the hypotheses of Theorem~\ref{four:thm:main}, the quantity \(A_e(S)\)
is finite and strictly positive for every end \(e\).
Truncating all other ends at sufficiently large coordinate spheres gives
a connected, complete one-ended exterior \(M_T\) with compact smooth
boundary \(S_T\). All its boundary components have future nonpositive
expansion with the normal into \(M_T\). Its selected-end charges are the
original charges, and
\begin{equation}\label{four:end:cut-truncation}
 A_*(M_T,g)\ge A_e(S),\qquad
 \lim_{T\to\infty}A_*(M_T,g)=A_e(S).
\end{equation}
The common parameter \(T\) in the limit can be replaced by any exhaustion
in which every unwanted truncation radius tends to infinity.
\end{lemma}

\begin{proof}
A sufficiently distant sphere in the chosen end bounds an admissible
outer domain, so the infimum is finite. To obtain a positive lower bound,
choose an embedded arc from an open disk in \(S\) to the chosen end, and
continue it along a straight coordinate ray. The compact part of the arc
can be chosen embedded and disjoint from the boundary except at its
initial point. A boundary collar and the tubular neighborhood theorem
give a proper tube with coordinates
\[
 [0,\infty)\times B^3\longrightarrow M.
\]
Its initial section lies in \(S\); its far part is a Euclidean tube of
fixed transverse coordinate size. Choose a smooth nonnegative function
\(\eta\) compactly supported in \(B^3\), with integral one, and define
in this tube
\[
 \alpha=\eta(z)\,\dd z^1\wedge\dd z^2\wedge\dd z^3.
\]
Extend \(\alpha\) by zero across the lateral tube boundary. It is a
smooth closed three-form on \(M\). Its comass, meaning the supremum of
its absolute value on unit simple three-vectors, has a finite bound
\(C_\alpha\): the initial part is compact, and the far metric is
uniformly comparable with the Euclidean metric. A sufficiently distant
selected-end sphere has \(\alpha\)-integral one, after fixing orientation.

Let \(D_e\) be any admissible outer domain with cut \(\Gamma\). Cut off
its selected-end tail at a sphere beyond \(\Gamma\). The resulting
domain is compact, since \(D_e\) contains no distant part of another
end. Its boundary is the distant sphere together with the \emph{entire}
intrinsic boundary \(\Gamma\), with the boundary orientations.
Stokes' theorem gives
\[
 \left|\int_\Gamma\alpha\right|=1,
 \qquad |\Gamma|_g\ge C_\alpha^{-1}.
\]
This proof applies to disconnected cuts and to frontier on \(S\).
In particular, it proves positivity before any end is truncated.

On an unwanted end, the normal of the truncation sphere that points
into the retained manifold is the inward radial normal. The given decay
therefore gives
\[
 H=-3T^{-1}+O(T^{-1-q}),\qquad
 \tr_{S_T}K=O(T^{-1-q}).
\]
The new boundary has strictly negative future expansion for large
\(T\). The other boundary components are the original ones.
Removing the open tails leaves a connected manifold: a path entering a
removed tail can have that portion replaced by a path along its connected
spherical cross-section. The retained manifold is closed in \(M\).
It is complete for its intrinsic metric as well. Indeed, an intrinsic
Cauchy sequence is ambient Cauchy and has a limit in the retained set;
near a smooth boundary or an interior point, the intrinsic and ambient
local distances are comparable in a half-ball or ball chart.

Every admissible outer domain in \(M_T\), viewed as a submanifold of
\(M\), is an admissible outer domain for \(e\). A part of its frontier
on a truncation sphere becomes a hypersurface in the interior of \(M\),
and remains part of its intrinsic boundary. This proves the first
inequality in \eqref{four:end:cut-truncation}. Conversely, a fixed original
admissible domain has compact frontier and contains no tail of any
unwanted end. It is consequently contained in \(M_T\) for all
sufficiently large truncation radii and is then an admissible domain
there. Applying this observation to cuts with area within any prescribed
positive error of \(A_e(S)\) proves the limit. The selected-end data have
not been changed, so neither selected-end charge changes.
\end{proof}

We may thus work on one-ended data until the final transfer. The boundary
is always included. No finite-perimeter relaxation or minimizing cut is
needed in this section.

\subsection{Conformal changes and an energy-raising shell}

\begin{lemma}[Conformal constraint and boundary formulas]\label{four:lem:conformal}
Let \(u>0\) be smooth, and set \(g_u=u^2g\), \(K_u=uK\). In dimension
four, as an identity of scalar functions and an identity of covectors,
respectively,
\begin{align}
 u^2\mu_u&=\mu-3u^{-1}\Delta_g u,\label{four:end:conformal-mu}\\
 uJ_u&=J+3K(\grad_g\log u,\cdot).\label{four:end:conformal-J}
\end{align}
Consequently,
\begin{equation}\label{four:end:conformal-dec}
 u^2\bigl(\mu_u-|J_u|_{g_u}\bigr)
 \ge \mu-|J|_g+3u^{-1}
          \bigl(-\Delta_g u-|K|_g|\dd u|_g\bigr).
\end{equation}
For any consistently oriented hypersurface,
\begin{equation}\label{four:end:conformal-expansion}
 \Theta_{K_u}=u^{-1}
       \bigl(\Theta_K+3\partial_\nu\log u\bigr).
\end{equation}
\end{lemma}

\begin{proof}
Put \(\varphi=\log u\). The scalar curvature formula and the tensor
scalings are
\[
 R_{g_u}=u^{-2}(R_g-6\Delta_g\varphi-6|\dd\varphi|_g^2),
 \qquad \tau_u=u^{-1}\tau,\qquad
 |K_u|_{g_u}^2=u^{-2}|K|_g^2.
\]
Since \(\Delta_g\varphi+|\dd\varphi|_g^2=u^{-1}\Delta_g u\), these
give \eqref{four:end:conformal-mu}. Write \(P=K-\tau g\), so that
\(P_u=uP\). The connection difference is
\[
 \widehat\Gamma^a_{ij}-\Gamma^a_{ij}
 =\delta_i^a\varphi_j+\delta_j^a\varphi_i-g_{ij}\varphi^a.
\]
Contracting the covariant derivative of \(uP\) gives
\[
 uJ_{u,i}=J_i+3P_{ij}\varphi^j-(\tr_gP)\varphi_i
          =J_i+3K_{ij}\varphi^j,
\]
because \(\tr_gP=-3\tau\). This displays the cancellation of all
trace-gradient terms. The covector norm scales by \(u^{-1}\), whence
\[
 u^2|J_u|_{g_u}=|J+3K(\grad\log u,\cdot)|_g.
\]
The triangle inequality proves \eqref{four:end:conformal-dec}.
Finally \(\nu_u=u^{-1}\nu\),
\(H_{g_u}=u^{-1}(H_g+3\partial_\nu\log u)\), and
\(\tr_{\mathrm{tan},g_u}K_u=u^{-1}\tr_{\mathrm{tan},g}K\), proving
\eqref{four:end:conformal-expansion}.
\end{proof}

\begin{lemma}[Positive shell]\label{four:end:positive-shell}
For any one-ended data in Theorem~\ref{four:thm:main} and every \(\lambda>0\),
there is a conformal change preserving all its hypotheses, with
\[
 E_\lambda=E+2\lambda,\qquad P_\lambda=P,
 \qquad A_*(g_\lambda)\ge A_*(g).
\]
The conformal factor is constant near the compact boundary and is at
least one everywhere.
\end{lemma}

\begin{proof}
Choose \(0<\xi<\min(q,1)\). For
\(T(r)=r^{-2}-r^{-2-\xi}\), at sufficiently large radius,
\[
 T'<0,\qquad
 -\Delta_gT-|K|_g|\dd T|_g
 =\xi(2+\xi)r^{-4-\xi}+O(r^{-4-q})>0.
\]
Take a nondecreasing smooth switch \(\beta\) from zero to one in this
region, constant on neighborhoods of its endpoints, and set
\[
 G(r)=-\int_r^\infty\beta(s)T'(s)\,\dd s.
\]
Extend \(G\) constantly throughout the interior. Then \(G\ge0\), it
equals \(T\) sufficiently far out, and
\[
 -\Delta_gG-|K|_g|\dd G|_g
 =\beta(-\Delta_gT-|K|_g|\dd T|_g)
       -\beta'T'|\dd r|_g^2\ge0.
\]
Use \(u=1+\lambda G\) in Lemma~\ref{four:lem:conformal}. The dominant
energy condition is preserved, as is the boundary expansion sign.
The new densities are integrable: outside a compact set the additional
terms in \eqref{four:end:conformal-mu}--\eqref{four:end:conformal-J} are bounded
by constants, depending on \(\lambda\), times
\(r^{-4-\xi}+r^{-4-q}\). The new falloff exponent can be taken to be
\(\min(q,2)>1\).

For completeness, the leading change in the metric flux numerator is
\(-6\partial_i(\lambda G)\). Products involving
\(g-\delta\), \(u-1\), or their first derivatives have vanishing
limiting flux. Hence \eqref{four:intro:energy} gives
\[
 E_\lambda-E=-\frac{\lambda}{\omega}
       \lim_{r\to\infty}\int_{S_r}\partial_rG\,\dd A_\delta
       =2\lambda.
\]
The transformed momentum tensor is exactly \(u(K-\tau g)\). Its
additional flux is \(O(r^{-q})\), so \(P_\lambda=P\) by
\eqref{four:intro:momentum}. Since \(u\ge1\), every cut has at least its
original area, proving the cut comparison.
\end{proof}

The shell will be used only after proving the inequality for
\(E>|P|\). It then excludes all remaining ADM vectors; no preliminary
strict timelikeness hypothesis will remain.

\subsection{A reference end with prescribed timelike charges}

\begin{lemma}[Explicit reference-plane charges]\label{four:end:boosted-plane}
For any \(E>|P|\), put \(m=(E^2-|P|^2)^{1/2}\). A spacelike plane
in the Schwarzschild--Tangherlini spacetime of mass \(m\), sufficiently
far from its central region, has induced vacuum data \((g_*,K_*)\)
with charges exactly \((E,P)\), and
\[
 g_*-\delta=O_k(r^{-2}),\qquad K_*=O_k(r^{-3})
 \quad\text{for every integer }k\ge0.
\]
Here \(K_*\) is one half the metric rate in the future normal direction.
\end{lemma}

\begin{proof}
The exterior isotropic form of the static vacuum metric
\cite{Tangherlini1963} is
\begin{equation}\label{four:end:tangherlini}
 \mathbf g=-\left(\frac{1-m/(2|y|^2)}{1+m/(2|y|^2)}\right)^2\dd b^2
       +\left(1+\frac{m}{2|y|^2}\right)^2\sum_{i=1}^4(\dd y^i)^2.
\end{equation}
For \(P=0\) use \(b=0\). For the computation in general, rotate the
first axis into the desired momentum direction, take \(0\le v<1\),
and write \(\gamma=(1-v^2)^{-1/2}\). Introduce inertial coordinates by
\[
 b=\gamma(t+vx_1),\qquad y_1=\gamma(x_1+vt),\qquad
 y_A=x_A\quad(A=2,3,4).
\]
The plane \(t=0\) is \(b=vy_1\). Its induced Minkowski metric is
\(\delta\); for sufficiently large radius it is spacelike for
\(\mathbf g\) as well. Let
\(s^2=\gamma^2x_1^2+|x_\perp|^2\). The leading spacetime perturbation
is \(m|y|^{-2}(2\dd b^2+\sum_i(\dd y^i)^2)\). On the plane its
nonzero components in the new coordinates are
\begin{equation}\label{four:end:boost-perturbation}
 \begin{aligned}
 h_{00}&=m(3\gamma^2-1)s^{-2},&
 h_{01}&=3m\gamma^2v s^{-2},\\
 h_{11}&=m(3\gamma^2-2)s^{-2},&
 h_{AB}&=m s^{-2}\delta_{AB},
 \end{aligned}
\end{equation}
with remainders \(O_k(r^{-4})\).
It follows by differentiating the spatial components that
\begin{equation}\label{four:end:boost-energy-integrand}
 \partial_j(g_*)_{ij}-\partial_i(g_*)_{jj}
       =6m\gamma^2 x_i s^{-4}+O(r^{-5}).
\end{equation}

To verify the momentum sign and normalization directly, the future
second form, to the order that contributes to the flux, is
\[
 (K_*)_{ij}=\tfrac12
       (\partial_t h_{ij}-\partial_i h_{0j}-\partial_j h_{0i})
       +O(r^{-5}).
\]
In addition to the spatial derivatives of \(s^{-2}\), use
\(\partial_t(s^{-2})=-2\gamma^2v x_1s^{-4}\) at \(t=0\).
One obtains
\[
 K_* =\frac{m\gamma^2v}{s^4}
 \begin{pmatrix}
 (3\gamma^2+2)x_1&3x_\perp^{\mathsf T}\\
 3x_\perp&-x_1 I_3
 \end{pmatrix}+O(r^{-5}),
\]
and therefore
\begin{equation}\label{four:end:boost-momentum-integrand}
 K_*-(\tr_{g_*}K_*)g_*
 =\frac{3m\gamma^2v}{s^4}
 \begin{pmatrix}
 x_1&x_\perp^{\mathsf T}\\
 x_\perp&-\gamma^2x_1 I_3
 \end{pmatrix}+O(r^{-5}).
\end{equation}
All discarded terms have vanishing sphere flux.

For \(n\in S^3\), let
\(D(n)=\gamma^2 n_1^2+|n_\perp|^2\). The ellipsoid
\(\gamma^2x_1^2+|x_\perp|^2\le1\) has volume
\(\omega/(4\gamma)\) by a linear change of variables. Polar integration
gives the same volume as \(\frac14\int_{S^3}D(n)^{-2}\,\dd A\).
Thus
\begin{equation}\label{four:end:ellipsoid-average}
 \int_{S^3}D(n)^{-2}\,\dd A=\omega/\gamma.
\end{equation}
Equations \eqref{four:intro:energy}, \eqref{four:end:boost-energy-integrand}, and
\eqref{four:end:ellipsoid-average} give \(E_*=m\gamma\).
The first row of \eqref{four:end:boost-momentum-integrand} gives
\(P_{*1}=m\gamma v\) with the denominator \(3\omega\) in
\eqref{four:intro:momentum}. Each transverse momentum integrand is a constant
multiple of \(n_1n_A D(n)^{-2}\), and integrates to zero by reflection.
Choose \(v=|P|/E\) and the rotation specified above.
The constraints are vacuum because these are induced hypersurface data
in the vacuum metric \eqref{four:end:tangherlini}. The displayed expansions
also give the asserted differentiated decay.
\end{proof}

\subsection{Compact inverses for the linear constraints}

On Euclidean four-space write
\begin{equation}\label{four:end:linear-operators}
 \mathcal L b=\partial_i\partial_j
                  (b_{ij}-(\tr_\delta b)\delta_{ij}),\qquad
 (\mathcal D k)_i=\partial_j
                  (k_{ij}-(\tr_\delta k)\delta_{ij}).
\end{equation}
These are the linearizations at \((\delta,0)\) of \(2\mu\) and \(J\).
The next lemma identifies every compatibility condition needed for a
compact correction.

\begin{lemma}[Compact linear constraint corrections]\label{four:end:compact-inverses}
Fix a compact subannulus \(U_0\) of a bounded connected Euclidean annulus
\(U\). If \(f\) and \(F\) are smooth and supported in \(U_0\), and
\begin{equation}\label{four:end:compatibility}
 \begin{aligned}
 \int f\phi\,\dd z&=0 &&\text{for every affine }\phi,\\
 \int F_iY^i\,\dd z&=0 &&\text{for every Euclidean Killing field }Y,
 \end{aligned}
\end{equation}
there are symmetric smooth tensors \(b,k\), supported in a fixed compact
subset of \(U\), with \(\mathcal Lb=f\), \(\mathcal Dk=F\).
For every integer \(j\ge0\) and every \(1<p_1<\infty\), they can be
chosen linearly in the sources so that
\begin{equation}\label{four:end:inverse-bounds}
 \|b\|_{W^{j+2,p_1}}\le C\|f\|_{W^{j,p_1}},\qquad
 \|k\|_{W^{j+1,p_1}}\le C\|F\|_{W^{j,p_1}}.
\end{equation}
The constant depends only on \(U_0,U,j,p_1\) and fixed localization
choices. The moment conditions are also necessary for compactly
supported solutions.
\end{lemma}

\begin{proof}
We first describe the scalar divergence inverse that will be used a
finite number of times. On a ball the regularized Bogovski\u\i\ operator
sends compactly supported smooth mean-zero data to compactly supported
smooth vector fields with the prescribed divergence, gaining one
\(W^{j,p_1}\) derivative. These are precisely the ball support and
Sobolev properties of the regularized operator
\cite{CostabelMcIntosh2010}.

To use it inside \(U\), cover a connected compact enlargement of the
source support by finitely many balls with closures in \(U\). Choose
a subordinate smooth partition and a spanning tree in the intersection
graph. For each non-root ball choose a smooth bump of integral one in
its overlap with its parent. Starting at the leaves, subtract from the
piece in that ball its integral times the overlap bump, and add that
multiple of the bump to its parent's piece. The corrected piece has
mean zero and remains compactly supported in its ball. At the root the
remaining integral is zero because the original source has integral
zero. Solve each piece in its ball and sum the resulting vector fields.
All transfers and estimates involve a fixed finite family of bumps,
so this gives a linear inverse with one derivative gain and support in
a fixed compact subset of \(U\). Choose successively larger compact
subsets in advance when several inverses will be used. This permits
every later inverse to act on the support produced by the preceding
one, without approaching the annular boundary.

For the scalar constraint, solve \(\partial_i v_i=f\). Compact support
and the linear moments give
\[
 \int v_i\,\dd z=-\int z_i f\,\dd z=0.
\]
Solve in turn \(\partial_jB_{ij}=v_i\), component by component.
The symmetric tensor \(S=(B+B^{\mathsf T})/2\) still satisfies
\(\partial_i\partial_jS_{ij}=f\), since the double divergence of a
skew tensor vanishes. The trace reversal in dimension four is inverted
by
\begin{equation}\label{four:end:trace-inverse}
 b=S-\tfrac13(\tr_\delta S)\delta.
\end{equation}
It gives \(b-(\tr_\delta b)\delta=S\), and hence \(\mathcal Lb=f\).
The two divergence inverses give the first estimate in
\eqref{four:end:inverse-bounds}.

For the vector constraint, the translational moments first allow a
matrix \(B\) with \(\partial_jB_{ij}=F_i\). Testing against the
rotational fields gives
\[
 \int(B_{ij}-B_{ji})\,\dd z=0\quad\text{for every }i,j.
\]
Apply the divergence inverse to obtain tensors \(D_{ijk}=-D_{jik}\)
such that
\[
 \partial_kD_{ijk}=-\tfrac12(B_{ij}-B_{ji}).
\]
Set
\[
 C_{ijk}=D_{ijk}-D_{ikj}-D_{jki}.
\]
Using only \(D_{ijk}=-D_{jik}\) gives the two identities
\[
 C_{ijk}=-C_{ikj},\qquad
 \tfrac12(C_{ijk}-C_{jik})=D_{ijk}.
\]
Consequently \(S_{ij}=B_{ij}+\partial_kC_{ijk}\) is symmetric and
\(\partial_jS_{ij}=F_i\): its skew part cancels, while the added row
divergence vanishes by the first identity. Apply
\eqref{four:end:trace-inverse} to \(S\) to obtain \(k\).
Here \(B\) has one derivative more than \(F\), \(D\) has two,
and the final differentiation loses one, proving the second estimate.
All operations preserve the stated support and smoothness.

Finally, integration by parts pairs \(\mathcal Lb\) with
\(\Hess\phi-(\Delta\phi)\delta\), which vanishes for affine
\(\phi\). It pairs \(\mathcal Dk\) with the symmetric gradient of
a Killing field, which also vanishes after trace reversal. This proves
necessity of \eqref{four:end:compatibility}.
\end{proof}

\subsection{Replacement, vacuum bending, and repair}

\begin{proposition}[Rest-end replacement]\label{four:prop:rest-end}
Let \((M,g,K)\) satisfy the one-ended hypotheses of
Theorem~\ref{four:thm:main}, and suppose \(E>|P|\). Put
\(m=(E^2-|P|^2)^{1/2}\). There is a sequence of smooth data
\((g_R,K_R)\) on the same manifold, complete with the same compact
boundary, satisfying the dominant energy condition and the same future
boundary inequality, such that:
\begin{enumerate}[label=\textup{(\roman*)}]
 \item \(K_R\) is compactly supported, and the ultimate end is an exact
 static Schwarzschild--Tangherlini spatial end;
 \item its energy is \(m_R=m+o(1)\), its momentum is zero, and both
 constraint densities are integrable;
 \item \(A_*(g_R)\ge(1-o(1))A_*(g)\).
\end{enumerate}
Every error here tends to zero as \(R\to\infty\) for the fixed original
data and fixed timelike vector. No uniformity as \(E-|P|\) tends to zero
is asserted or needed.
\end{proposition}

\begin{proof}
Use Lemma~\ref{four:end:boosted-plane} for a reference end with exactly the
original charges. We first join the data to this plane on
\(R<r<4R\), controlling the full pointwise constraint deficit.

\paragraph{The scaled commutators.}
Fix \(1<p<\min(2,q)\), write \(x=Rz\), and on
\(1<|z|<4\) use the scaled component fields
\[
 g(Rz),\qquad k(z)=RK(Rz),
\]
and their reference counterparts. This corresponds to rescaling lengths
by \(R^{-1}\), so both constraint densities scale by \(R^2\).
The metric deviations and tensors have size \(O(R^{-p})\) in
\(C^2\) and \(C^1\), respectively. Expanding the constraint map about
\((\delta,0)\) gives its linear part
\((\mathcal L,\mathcal D)\) from \eqref{four:end:linear-operators}, with
remainder bounded in \(C^0\) by \(CR^{-2p}\). This follows directly
from the scalar curvature formula: the nonlinear metric terms are
bounded by \(C(|b||D^2b|+|Db|^2)\); the other scalar terms are quadratic
in \(k\). The momentum remainder is bounded by
\(C(|b||Dk|+|Db||k|)\).

Choose a fixed smooth radial \(\psi\), equal to one near \(|z|=1\)
and zero near \(|z|=4\), with \(0\le\psi\le1\). Blend the metric
components and tensors by \(\psi\). If \(b\) is the difference of
the two scaled metric perturbations and \(k\) the difference of the
two scaled tensors, the errors of the linear constraints relative to
the blended sources are
\(f_R=[\mathcal L,\psi]b\) and
\(F_R=[\mathcal D,\psi]k\). For \(T_{ij}=b_{ij}-(\tr_\delta b)\delta_{ij}\),
\begin{align}
 [\mathcal L,\psi]b
   &=(\partial_i\partial_j\psi)T_{ij}
       +2(\partial_i\psi)\partial_jT_{ij},\label{four:end:scalar-commutator}\\
 ([\mathcal D,\psi]k)_i
   &=(\partial_j\psi)(k_{ij}-(\tr_\delta k)\delta_{ij}).
       \label{four:end:vector-commutator}
\end{align}
These expressions are supported in a fixed compact subannulus, and
their \(C^1\) norms are \(O(R^{-p})\). In particular, the first bound
uses only two metric derivatives, and the second only one tensor
derivative.

\paragraph{All compatibility moments are small.}
We prove the sharper bound \(o(R^{-2})\) for every moment in
\eqref{four:end:compatibility}. Work momentarily in the physical coordinates,
with
\[
 h=g-g_*,\quad
 T=K-(\tr_\delta K)\delta-K_*+(\tr_\delta K_*)\delta,
 \quad f=\mathcal Lh,\quad F_i=\partial_jT_{ij}.
\]
Both \(f\) and \(F\) belong to \(L^1\) on the end. Indeed, the
difference between the physical constraints and their Euclidean linear
parts is \(O(r^{-2-2\min(q,2)})\), which is integrable in dimension
four. The physical densities are integrable by hypothesis and by
reference vacuum, and Euclidean and physical measures and norms are
uniformly comparable there.

An elementary consequence of integrability is
\begin{equation}\label{four:end:weak-first-moment}
 \int_{r_0<r<4R}r(|f|+|F|)\,\dd x=o(R).
\end{equation}
To see this, split at a fixed radius \(L\). The integral up to \(L\),
divided by \(R\), tends to zero. The remaining integral divided by
\(R\) is at most four times the \(L^1\) tail beyond \(L\), which can
be made arbitrarily small. A finite first moment is not required.

For an affine function \(\phi\), define the flux primitive
\[
 Q_{\phi,i}
 =\phi(\partial_jh_{ij}-\partial_i h_{jj})
       -(\partial_j\phi)h_{ij}+(\partial_i\phi)h_{jj}.
\]
Its divergence is \(\phi f\). For a Killing field \(Y\), symmetry
of \(T\) gives \(\partial_j(Y^iT_{ij})=Y^iF_i\).
Let \(\mathfrak F_\phi(r)\) and \(\mathfrak G_Y(r)\) be the corresponding
outward Euclidean sphere fluxes. For constant tests, matching the ADM
charges gives
\[
 \mathfrak F_1(r)=o(1),\qquad
 \mathfrak G_{e_i}(r)=o(1).
\]
In the momentum flux the replacement of \(\tr_gK\,g\) by
\((\tr_\delta K)\delta\) costs at most
\(O(r^{2-2\min(q,2)})=o(1)\). For homogeneous linear scalar tests and
rotational fields, integration from a fixed sphere and
\eqref{four:end:weak-first-moment} instead give
\[
 \mathfrak F_\phi(r)=o(r),\qquad \mathfrak G_Y(r)=o(r).
\]

Put \(\psi_R(x)=\psi(x/R)\), and let \(A_R=\{R<r<4R\}\).
Integration by parts, with \(\psi_R=1\) near the inner sphere and zero
near the outer sphere, gives the exact formulas
\begin{align}
 \int_{A_R}\phi[\mathcal L,\psi_R]h\,\dd x
  &=-\mathfrak F_\phi(R)-\int_{A_R}\psi_R\phi f\,\dd x,
       \label{four:end:scalar-moment}\\
 \int_{A_R}Y^i([\mathcal D,\psi_R](K-K_*))_i\,\dd x
  &=-\mathfrak G_Y(R)-\int_{A_R}\psi_RY^iF_i\,\dd x.
       \label{four:end:vector-moment}
\end{align}
In scaled coordinates, constant moments have the factor \(R^{-2}\):
the source scales by \(R^2\) and volume by \(R^{-4}\).
A linear test in \(z\) contributes the additional factor \(R^{-1}\).
Thus \eqref{four:end:scalar-moment}--\eqref{four:end:vector-moment} and the preceding
estimates show
\begin{equation}\label{four:end:small-moments}
 \int f_R\phi\,\dd z=o(R^{-2}),\qquad
 \int (F_R)_iY^i\,\dd z=o(R^{-2})
\end{equation}
for each element of a fixed basis of the affine functions or Killing
fields.

\paragraph{Correction and the physical deficit.}
Choose a nonnegative smooth bump \(\zeta\), supported in a fixed ball
within the annulus and positive on a smaller ball. For a basis of the
affine tests, its Gram matrix
\(\int\zeta\phi_a\phi_b\) is positive definite. For a basis of
Killing fields, the matrix \(\int\zeta\,Y_a\cdot Y_b\) is likewise
positive definite: a nonzero affine Killing field cannot vanish on an
open ball. Subtract the corresponding bump combinations from
\(f_R,F_R\) to erase their moments. By \eqref{four:end:small-moments},
the coefficients, and every fixed smooth norm of these bump corrections,
are \(o(R^{-2})\).

Apply Lemma~\ref{four:end:compact-inverses} to the negatives of the
moment-free commutators and add the resulting tensors to the blend.
With \(j=1\) and \(p_1>4\), Sobolev embedding and
\eqref{four:end:inverse-bounds} bound the metric correction in \(C^2\)
and the tensor correction in \(C^1\) by \(CR^{-p}\).
Their supports stay in a fixed enlarged compact subannulus.
The new metric is positive for sufficiently large \(R\) and agrees
exactly with the original data inside \(R\) and the plane outside
\(4R\).

Let \((\widetilde g_R,\widetilde K_R)\) denote these physical data.
Their scaled constraints equal \(\psi\) times the original scaled
constraints, plus an error
\(o(R^{-2})+O(R^{-2p})\); the reference constraints vanish.
The change in the momentum norm costs another \(O(R^{-2p})\), because
the metric change is \(O(R^{-p})\) and the scaled momentum density is
\(O(R^{-p})\). Since \(\psi\ge0\), the original dominant energy
condition implies, after undoing the scaling,
\begin{equation}\label{four:end:physical-deficit}
 \widetilde\mu_R-|\widetilde J_R|_{\widetilde g_R}
       \ge-\eta_RR^{-4},\qquad \eta_R\longrightarrow0.
\end{equation}
Indeed the remaining physical error is
\(R^{-2}[o(R^{-2})+O(R^{-2p})]=o(R^{-4})\), since \(p>1\).
The possible deficit is confined to a fixed closed subannulus on the
scale \(R\). Replace \(\eta_R\) by a positive majorant tending to zero
if necessary.

\paragraph{Bending the plane to a static slice.}
Only the exact reference region is used for this step. In its spatial
coordinates \(y\), write the plane as \(b=v\cdot y\). Choose a smooth
switch \(\psi_1\) from one to zero on \([0,1]\), constant outside that
interval, and replace the plane farther out by the graph
\begin{equation}\label{four:end:vacuum-bend}
 b=(v\cdot y)\psi_1\left(L^{-1}\log(|y|/R_b)\right).
\end{equation}
Here \(R_b\) is a sufficiently large fixed multiple of \(R\), so the
graph initially lies wholly beyond the gluing annulus. For fixed
\(|v|<1\), choose the fixed \(L\) so large that
\[
 |\grad_y b|
 \le |v|\bigl(1+\|\psi_1'\|_\infty/L\bigr)<1.
\]
The Minkowski graph is uniformly spacelike. Since the perturbation
\eqref{four:end:tangherlini} is \(O(r^{-2})\), it remains uniformly spacelike
in that metric for large \(R\). It agrees with the plane on an open
inner region and with \(b=0\) on an open outer region. These induced
data are vacuum everywhere in the bend, so it adds no deficit.

In the fixed \(x\) chart, use \(y=Ax\) where
\(A=(I-v\otimes v)^{-1/2}\), with the chosen rotation. The derivatives
of \eqref{four:end:vacuum-bend} give \(D^2b=O(r^{-1})\), and the induced
data, now denoted \((g_R^0,K_R^0)\), satisfy on all transition regions
\begin{equation}\label{four:end:transition-geometry}
 c\delta\le g_R^0\le C\delta,\qquad
 |\partial g_R^0|+|K_R^0|\le C/r.
\end{equation}
Here \(c,C>0\) may depend on the fixed boost and switches, but not
on large \(R\). Outside a fixed multiple of \(R\), the data are static
in the \(y\) chart and have tensor zero.

For the repair choose a smooth radius \(\mathfrak r\), equal to
\(|x|\) through the gluing annulus and through the bend, and equal to
\(|y|\) farther out in the static region. It can be chosen with
\begin{equation}\label{four:end:radius-bounds}
 \mathfrak r\asymp r,\qquad
 c\le|\dd\mathfrak r|_{g_R^0}\le C,
 \qquad |\Delta_{g_R^0}\mathfrak r|\le C/r.
\end{equation}
Here is an explicit way to make the interpolation. In the static region
put \(a(n)=|A^{-1}n|\), \(n=y/|y|\), and interpolate
\[
 \log\mathfrak r=\log|y|+(1-\chi(\log(|y|/R_c)))\log a(n).
\]
Take \(R_c\) beyond the bend and make \(\chi\) change from zero to one
over a sufficiently long fixed logarithmic interval. Its slope can be
made so small that the derivative of \(\log\mathfrak r\) with respect
to \(\log|y|\) lies between \(1/2\) and \(3/2\). Angular derivatives
are bounded, second spatial derivatives are \(O(r^{-1})\), and
\eqref{four:end:transition-geometry} proves \eqref{four:end:radius-bounds}.
All transition intervals are contained between fixed multiples of \(R\).

\paragraph{A conformal repair with vanishing mass cost.}
Write \(s=\mathfrak r/R\). We construct a nonnegative function
\(F_R(s)\), constant on the inside, and use
\[
 u_R=1+\varepsilon_RR^{-2}F_R(s),\qquad
 \varepsilon_R=C_0\eta_R.
\]
Choose a fixed interval \([a,b_1]\) in the \(s\) variable that starts
before the possible deficit, contains all transition regions, and ends
where the radius is \(|y|\) and the data are static. This interval lies
in the end for all large \(R\). Put \(z_R=-F_R'\) and prescribe
\begin{equation}\label{four:end:repair-ode}
 z_R(a)=0,\qquad z_R'=C_1z_R+\omega_1,
\end{equation}
where \(\omega_1\ge0\) is smooth, vanishes on a neighborhood of \(a\),
and is at least one on a neighborhood of the possible deficit interval.
Choose it to vanish near \(b_1\). The solution is nonnegative and its
norms on this fixed interval are bounded independently of large \(R\).

Differentiating \(u_R\) and using \eqref{four:end:transition-geometry} and
\eqref{four:end:radius-bounds} gives fixed constants \(c',C'>0\) such that
\begin{equation}\label{four:end:repair-estimate}
 -\Delta_{g_R^0}u_R-|K_R^0||\dd u_R|_{g_R^0}
 \ge\varepsilon_RR^{-4}(c'z_R'-C'z_R)
 \quad\text{on }a\le s\le b_1.
\end{equation}
In fact the first positive term is
\(\varepsilon_RR^{-4}z_R'|\dd\mathfrak r|^2\); the other terms are
\(\varepsilon_RR^{-3}z_R\Delta\mathfrak r\) and the negative tensor
term, both bounded below by \(-C\varepsilon_RR^{-4}z_R\) here.
Choose \(C_1>C'/c'\). Equation \eqref{four:end:repair-ode} then makes
\eqref{four:end:repair-estimate} nonnegative everywhere on this interval
and at least \(c'\varepsilon_RR^{-4}\) on the deficit region.

On the static region put \(c_m=m/2\) and continue \(z_R\) so that
\[
 q_R(s)=s^3\left(1+\frac{c_m}{R^2s^2}\right)^2z_R(s)
\]
is nondecreasing and becomes a positive constant \(q_{R,\infty}\).
This can be done smoothly with a uniform bound on \(q_{R,\infty}\).
Indeed, immediately past \(b_1\), continue \eqref{four:end:repair-ode};
both its derivative and that of the positive prefactor make \(q_R'\)
nonnegative for large \(R\). Multiply this nonnegative derivative by
a smooth switch that equals one initially and zero after one more fixed
interval, and integrate. This preserves all matching derivatives and
keeps \(q_R\) nondecreasing. The radial Laplacian of the static metric
then shows \(-\Delta_{g_R^0}u_R\ge0\) there, since its radial divergence
is proportional to \(-q_R'\); also \(K_R^0=0\).

Set \(F_R(s)=\int_s^\infty z_R(t)\,\dd t\), extending it constantly
before \(a\). Its norm is bounded uniformly: the finite intervals have
bounded length and data, while on the tail \(z_R=O(s^{-3})\).
Thus \(1\le u_R\le2\) for large \(R\).
Choose the fixed \(C_0\) large enough that
\(3u_R^{-1}c'\varepsilon_R\ge\eta_R\). Lemma~\ref{four:lem:conformal},
\eqref{four:end:physical-deficit}, and \eqref{four:end:repair-estimate} now prove
the dominant energy condition everywhere for
\[
 g_R=u_R^2g_R^0,\qquad K_R=u_RK_R^0.
\]
The repair is constant in a neighborhood of the original boundary, and
therefore preserves its expansion sign.

On the ultimate tail the integral for \(F_R\) is explicit:
\[
 u_R=1+\frac{a_R}{|y|^2+c_m},\qquad
 a_R=\tfrac12\varepsilon_Rq_{R,\infty}=O(\varepsilon_R)=o(1).
\]
Consequently
\begin{equation}\label{four:end:static-repaired-metric}
 g_R=\left(1+\frac{c_m+a_R}{|y|^2}\right)^2\delta,
 \qquad K_R=0
 \quad\text{on that tail}.
\end{equation}
Its energy is \(m+2a_R=m+o(1)\), with zero momentum, by the direct
static instance of Lemma~\ref{four:end:boosted-plane}. The densities vanish
on this tail and are smooth on the remaining compact part, so they are
integrable. The data are complete with boundary: a fixed compact core
is smooth, and the end has a uniform positive metric lower bound for
each \(R\).

\paragraph{Comparison of every enclosing cut.}
The final metric dominates the original \(g\) for \(r\le R\), because
the intermediate data there are the original ones and \(u_R\ge1\).
For \(r\ge R\), \eqref{four:end:transition-geometry} and the static tail
give \(g_R\ge c\delta\) in the original \(x\) chart with a fixed
\(c>0\). Choose \(0<\kappa<\min(1,\sqrt c/2)\).
Let \(L_R=R^{1/2}\), and choose a smooth strictly increasing radial map
\(\rho_R\) equal to \(r\) for \(r\le L_R\), with derivative decreasing
from one to \(\kappa\) on \([L_R,2L_R]\), and derivative \(\kappa\)
afterward. It defines a diffeomorphism
\[
 \Phi_R(x)=\rho_R(|x|)\,\frac{x}{|x|}
\]
on the end, extended by the identity on the core. It is proper because
\(\rho_R(r)\to\infty\), preserves the selected end, and fixes the
boundary. Its radial and tangential Euclidean dilations are
\(\rho_R'\) and \(\rho_R/r\). They are at most one everywhere, and
for \(r\ge R\) are at most \(\kappa+O(R^{-1/2})\).

On the nonidentity region with \(r\le R\), both \(r\) and
\(\rho_R(r)\) tend uniformly to infinity; the original metric is
therefore \((1+o(1))\delta\) at both points. The dilation bound and
\(g_R\ge g\) show
\(\Phi_R^*g\le(1+o(1))g_R\) there. For \(r\ge R\), the choice
of \(\kappa\), the same original-end decay, and \(g_R\ge c\delta\)
give this inequality as well. It is immediate on the identity core.
We have thus proved the global tensor comparison
\begin{equation}\label{four:end:compression-metric}
 \Phi_R^*g\le(1+\epsilon_R)g_R,\qquad\epsilon_R\longrightarrow0.
\end{equation}
An admissible outer domain for \(g_R\) maps under this proper
diffeomorphism to an admissible original outer domain, with all
coincident frontier retained. Its three-dimensional area is at most
\((1+\epsilon_R)^{3/2}\) times the original cut's \(g_R\)-area.
Taking infima gives
\[
 A_*(g)\le(1+\epsilon_R)^{3/2}A_*(g_R),
\]
which proves the last assertion without a compactness assumption on
minimizing sequences.
\end{proof}

\subsection{Strictification with a controlled asymptotic tail}

\begin{lemma}[Strictification]\label{four:lem:strictification}
Let \((M,g,K)\) be smooth one-ended data with compact boundary,
complete with the boundary included, satisfying the dominant energy
condition and \(\Theta_K\le0\). Suppose \(K\) is compactly supported
and the ultimate end is an exact static Schwarzschild--Tangherlini
spatial end. Fix \(0<\delta_1<1\), and extend its radial coordinate to
a positive smooth function \(r\) on \(M\).
There are smooth conformal data
\((g_s,K_s)=(e^{2s\phi}g,e^{s\phi}K)\), for all sufficiently small
\(s>0\), such that
\begin{equation}\label{four:end:strict-data}
 \Theta_{K_s}<0,\qquad
 \mu_s-|J_s|_{g_s}\ge c_s r^{-4-\delta_1},\qquad c_s>0.
\end{equation}
They are complete, have compactly supported tensor, integrable constraint
densities, finite energy, and on the end
\begin{equation}\label{four:end:strict-decay}
 g_s-\delta=O_k(r^{-2})\quad(k\ge0),\qquad
 R_{g_s}=O(r^{-4-\delta_1}).
\end{equation}
Moreover \(E_{g_s}\to E_g\) and \(A_*(g_s)\to A_*(g)\) as
\(s\downarrow0\).
\end{lemma}

\begin{proof}
Choose a smooth compactly supported \(b_0\ge|K|_g\), and solve
\begin{equation}\label{four:end:strictification-equation}
 \begin{split}
 -\Delta_g\phi
   &=b_0\sqrt{|\dd\phi|_g^2+r^{-6}}+r^{-4-\delta_1},\\
 \partial_\nu\phi&=-1\quad\text{on }\partial M,
 \qquad\phi\longrightarrow0\quad\text{at infinity}.
 \end{split}
\end{equation}
The boundary normal here points into \(M\). We give the solvability
argument, including the estimates needed for the condition at infinity.

First truncate at an outer sphere \(S_T\) in the static region and
prescribe \(\phi=0\) there. Multiply \(b_0\) by a parameter
\(t_0\in[0,1]\). At zero the problem is the ordinary mixed Poisson
problem, with a nonempty Dirichlet face and disjoint smooth Neumann
faces. At any solution, the linearization for homogeneous variations is
\[
 v\longmapsto-\Delta_gv
    -t_0b_0\frac{\inner{\dd\phi}{\dd v}_g}
                    {\sqrt{|\dd\phi|_g^2+r^{-6}}},
 \qquad \partial_\nu v=0\text{ on }\partial M,\quad v=0\text{ on }S_T.
\]
The drift has magnitude at most \(b_0\). The maximum principle and
boundary point principle give a zero kernel: a nonzero positive maximum
or negative minimum cannot occur in the interior or on a homogeneous
Neumann face, and the Dirichlet values are zero. The mixed Laplacian
has Fredholm index zero; adding the smooth first-order term preserves
that index. Hence the linearization is invertible in the usual mixed
H\"older spaces. The linear estimates and maximum principles used here
are the ordinary scalar elliptic ones; see \cite{GilbargTrudinger2001}.

Fix a truncation and \(p_1>4\). The mixed Laplace estimate,
first-derivative interpolation, and the elementary bound
\[\sqrt{|\dd\phi|^2+r^{-6}}\le|\dd\phi|+r^{-3}\]
give
\begin{equation}\label{four:end:strictification-local-bound}
 \|\phi\|_{W^{2,p_1}(M_T)}
       \le C_T(1+\|\phi\|_{L^\infty(M_T)}),
\end{equation}
uniformly in \(t_0\). Nonnegativity follows from the minimum principle:
the equation has strictly positive right side; at an inner boundary
minimum the derivative into the domain cannot equal \(-1\), and the
outer boundary value is zero.

If the suprema on a fixed truncation were unbounded, divide a sequence
by its supremum. By \eqref{four:end:strictification-local-bound} and compact
Sobolev embedding, a subsequence converges in \(C^{1,\alpha}\) for
some \(\alpha>0\) to a nonnegative function \(v\) with supremum one,
homogeneous mixed data, and
\[
 -\Delta_gv=t_*b_0|\dd v|_g.
\]
The latter equation is a linear equation with bounded measurable drift
along \(v\): set the drift to
\(t_*b_0\grad v/|\dd v|\) where the gradient is nonzero, and to zero
elsewhere. The strong maximum and boundary point principles again
force \(v=0\), a contradiction. Thus the suprema are bounded.
The nonlinearity in \eqref{four:end:strictification-equation} is smooth on
a fixed truncation, since \(r^{-6}\) has a positive minimum there.
Sobolev embedding, Schauder estimates, and iteration now give all
required classical bounds. Invertibility gives openness and these
bounds give closedness of the parameter set. A solution therefore
exists at \(t_0=1\) on every sufficiently large truncation.

We next obtain bounds independent of \(T\). Fix \(r_0\) beyond the
support of \(b_0\) and in the static region. For large fixed \(r_0\),
the positive function
\[
 W(r)=r^{-2}(1-r^{-\delta_1})
\]
satisfies \(-\Delta_gW\ge c r^{-4-\delta_1}\) there, for some
\(c>0\). Its Euclidean leading coefficient is
\(\delta_1(2+\delta_1)\), while the static metric error is
\(O(r^{-6})\), dominated since \(\delta_1<1\).
Comparison with a sufficiently large multiple of
\((1+\sup_{r\le r_0}\phi)W\) gives
\begin{equation}\label{four:end:strictification-end-barrier}
 0\le\phi\le C(1+\sup_{r\le r_0}\phi)r^{-2}
       \quad(r_0\le r\le T),
\end{equation}
with \(C\) independent of \(T\). The comparison has the required
ordering on \(S_{r_0}\) by the chosen multiple and on \(S_T\) because
the solution vanishes there.

If the core suprema diverged along a sequence \(T\to\infty\), normalize
by those suprema. The local version of
\eqref{four:end:strictification-local-bound}, including the fixed inner
boundary, gives a subsequential limit with core supremum one. It solves
the same homogeneous bounded-drift equation as above, has homogeneous
inner Neumann data, and tends to zero at infinity by
\eqref{four:end:strictification-end-barrier}. Its positive global maximum
is attained in a compact set. The strong maximum and boundary point
principles exclude such a maximum. This contradiction gives uniform
core bounds. Local compactness and exhaustion now solve
\eqref{four:end:strictification-equation} on \(M\), with \(\phi\ge0\).

On the static tail the equation is simply
\(-\Delta_g\phi=r^{-4-\delta_1}\). Bound
\eqref{four:end:strictification-end-barrier} and rescaled Poisson estimates
on annuli imply
\begin{equation}\label{four:end:phi-decay}
 \phi=O_k(r^{-2})\quad\text{for every }k\ge0.
\end{equation}
The right side and static coefficients have differentiated decay of
all orders, so these estimates follow successively from the rescaled
\(W^{2,p_1}\) and Schauder estimates. The limit
\[
 L_\phi=\lim_{r\to\infty}\int_{S_r}\partial_r\phi\,\dd A_\delta
\]
exists: integration of \(\Delta_g\phi\), which is integrable, gives
the limit with physical normal and area; \eqref{four:end:phi-decay} and
\(g-\delta=O_1(r^{-2})\) make their difference tend to zero.

Now put \(u=e^{s\phi}\). Lemma~\ref{four:lem:conformal} gives
\begin{align*}
 u^2(\mu_s-|J_s|_{g_s})
 &\ge\mu-|J|+3s\bigl(-\Delta_g\phi
                       -|K||\dd\phi|-s|\dd\phi|^2\bigr)\\
 &\ge3s\bigl(r^{-4-\delta_1}-s|\dd\phi|^2\bigr).
\end{align*}
The function \(|\dd\phi|^2r^{4+\delta_1}\) is bounded: it is smooth
on the compact part and is \(O(r^{-2+\delta_1})\) on the end.
For all sufficiently small \(s>0\), the last expression is at least
\(\frac32s r^{-4-\delta_1}\). Boundedness of \(\phi\) then proves
the gap in \eqref{four:end:strict-data}. The boundary formula is
\[
 \Theta_{K_s}=e^{-s\phi}(\Theta_K-3s)<0.
\]
Equation \eqref{four:end:phi-decay} proves the metric falloff, and the
scalar conformal formula on the static tail gives
\(R_{g_s}=O(r^{-4-\delta_1})\). There \(K_s=0\), so both density
integrability and \eqref{four:end:strict-decay} follow. Completeness follows
from the global comparison \(g_s\ge g\).

Finally the same flux calculation as in Lemma~\ref{four:end:positive-shell}
gives the exact energy change
\[
 E_{g_s}-E_g=-\frac{s}{\omega}L_\phi.
\]
The remainder from the exponential and metric products has zero
limiting flux by \eqref{four:end:phi-decay}. Thus the energy is finite and
converges to \(E_g\). Pointwise,
\(g\le g_s\le e^{2s\|\phi\|_\infty}g\), so
\[
 A_*(g)\le A_*(g_s)
       \le e^{3s\|\phi\|_\infty}A_*(g).
\]
This proves the asserted convergence of the exact cut infima.
\end{proof}

\subsection{Transfer of the prepared energy estimate}

\begin{proposition}[Closure of the end reduction]\label{four:end:transfer}
The prepared energy estimate \eqref{four:end:prepared-interface} implies
Theorem~\ref{four:thm:main}.
\end{proposition}

\begin{proof}
First take one-ended original data with \(E>|P|\). For each sufficiently
large fixed \(R\), Proposition~\ref{four:prop:rest-end} supplies a repaired
rest end \((g_R,K_R)\) of energy \(m_R\to m\), where
\(m=(E^2-|P|^2)^{1/2}\). Apply Lemma~\ref{four:lem:strictification} to this
fixed pair. Its output has all the hypotheses of
Theorem~\ref{four:thm:prepared-energy}: a smooth connected orientable complete
one-ended exterior with nonempty compact smooth boundary, compactly
supported \(K\), strictly negative future boundary expansion, a positive
gap bounded below by a multiple of \(r^{-4-\delta_1}\), and the end
conditions \eqref{four:end:strict-decay}. That theorem gives
\[
 E_{g_{R,s}}\ge\frac12
       \left(\frac{A_*(g_{R,s})}{\omega}\right)^{2/3}.
\]
Let \(s\downarrow0\) at fixed \(R\). The strictification limits yield
\[
 m_R\ge\frac12\left(\frac{A_*(g_R)}{\omega}\right)^{2/3}.
\]
Now let \(R\to\infty\) and use the proper-map comparison in
Proposition~\ref{four:prop:rest-end}. The result is
\begin{equation}\label{four:end:timelike-conclusion}
 \sqrt{E^2-|P|^2}\ge
       \frac12\left(\frac{A_*(g)}{\omega}\right)^{2/3}
       \quad\text{when }E>|P|.
\end{equation}

Suppose next that the original one-ended charges satisfy \(E\le|P|\).
Choose any sequence \(a_j>0\) decreasing to zero and put
\(\lambda_j=(|P|+a_j-E)/2>0\). By
Lemma~\ref{four:end:positive-shell}, the conformally changed data have
energy \(|P|+a_j\), momentum \(P\), and cut infimum at least the
original positive \(A_*(g)\). They are timelike, so
\eqref{four:end:timelike-conclusion} applied separately to each gives
\[
 \sqrt{(|P|+a_j)^2-|P|^2}
       \ge\frac12\left(\frac{A_*(g)}{\omega}\right)^{2/3}>0.
\]
The left side tends to zero, a contradiction. Thus all original
one-ended data in this class have \(E>|P|\), and
\eqref{four:end:timelike-conclusion} holds for all of them. The argument
includes the cases of zero or negative energy and the null case.
It takes no limit of the replacement geometries as the boost degenerates.

For finite-ended original data, apply the complete one-ended conclusion
to a sufficiently distant truncation of every other end from
Lemma~\ref{four:lem:cut-reduction}. The charges at the selected end are
unchanged, and its one-ended cut infimum is at least \(A_e(S)\).
Therefore
\[
 E_e\ge\sqrt{|P_e|^2+
                 \frac14\left(\frac{A_e(S)}{\omega}\right)^{4/3}}.
\]
This is Theorem~\ref{four:thm:main} in the stated normalization. The only
remaining task is to prove \eqref{four:end:prepared-interface}; the following
sections do so by constructing a Riemannian metric and applying the
minimal-enclosure comparison in Section~\ref{four:sec:existence}.
\end{proof}
